%! TEX program = lualatex
\documentclass[14pt]{extarticle}

\usepackage[ukrainian]{babel}
\usepackage[margin=2cm]{geometry}

\usepackage{fontspec}
\setmainfont{CMU Serif}
\setsansfont{CMU Sans Serif}
\defaultfontfeatures{Ligatures={TeX}}

\usepackage{csquotes}

\usepackage[backend=biber, style=authoryear, sorting=nyt]{biblatex}
\addbibresource{references.bib}

\emergencystretch=1em
\hbadness=10000
\pagestyle{empty}

\begin{document}

\section*{Огляд літератури (biber, стиль authoryear)}

Загальні принципи адитивних технологій описано у~\parencite{Gibson2021}.
Проектування FFF-принтерів розглянуто у~\parencite{Soni2021},
а порівняльний аналіз картезіанських та дельта-принтерів~--- у~\parencite{Schmitt2018}.
Типові дефекти FDM-друку систематизовано у~\parencite{Baechle2022}.
Методи виявлення помилок на основі комп'ютерного зору досліджено
у~\parencite{Baumann2016, Langeland2020, Paulsen2019}.
Машинне навчання для контролю якості поверхні застосовано у~\parencite{Kadam2021}.
Методи оцінки якості зображень описано у~\parencite{Crete2007, Wang2005}.
Основи цифрової обробки зображень викладено у~\parencite{Gonzalez2018}.
Технічні характеристики комерційних 3D-принтерів наведено у~\parencite{BambuLabA1Mini}.
Моделювання вбудованого 3D-друку розглянуто у~\parencite{VanDriel2025}.

\nocite{*}

\printbibliography[title={Список використаних джерел}]

\end{document}
